% ============================================================
%  Глава 4. Данные и методы их обработки
% ============================================================
\chapter{Данные и методы обработки}
\label{ch:data}

\section{Цены базовых активов}
\label{sec:data-assets}

Дневные цены закрытия базовых активов выгружены через открытое
API Yahoo Finance. По ним посчитаны лог-доходности
\[
  r_t \;=\; \ln\bigl(S_t/S_{t-1}\bigr),
\]
а реализованная корреляция оценивается на скользящем окне в
$21$~рабочий день по матрице лог-доходностей корзины. Список тикеров и
конкретный временной диапазон, использованные в эксперименте, приведены
в~\cref{app:data-details}.

\todoblock{Указать диапазон дат (например, 2003-01-01 --- 2025-12-31)
после фиксации финального датасета.}

\section{Цены корреляционных свопов}
\label{sec:data-corrswap}

Цены корреляционных свопов в обучающем датасете получены через
приближённые \emph{замкнутые} формулы из~\cref{sec:corr-models}:
константная модель, модель возврата к среднему, модель Якоби и
приближение Фишера--OU. Использование закрытых формул вместо Монте-Карло
оправдано двумя соображениями:

\begin{enumerate}
  \item Замкнутые формулы дают разумное приближение в реалистичных
        диапазонах параметров $(\kappa,\bar\rho,\sigma,\eta)$.
  \item Монте-Карло на полном пласте параметров (десятки тысяч точек
        $\times$ десятки моделей) был бы вычислительно невозможен.
\end{enumerate}

Параметры моделей $(\kappa,\bar\rho,\sigma)$ оцениваются по годовому
окну исторических реализованных корреляций методом максимального
правдоподобия (для модели Якоби) и МНК (для остальных).

\placeholderfig{12cm}{6cm}%
{Распределение оцененных параметров $(\kappa,\bar\rho)$ по годовым
окнам.}{fig:data-params}

\section{Цены rainbow-опционов}
\label{sec:data-rainbow}

Для rainbow-опционов референсные цены получены численно
\todoblock{Указать конкретный метод (двумерная PDE, специализированное
квазимонте-карло, формула Stulz/Johnson) после фиксации.}, а не
полноценным Монте-Карло: при объёме обучающей выборки в
$\sim 10^{5}\!-\!10^{6}$ точек прямой Монте-Карло на каждую был бы
вычислительно неприемлем.

\section{Дополнительная предобработка}
\label{sec:data-prep}

Перед обучением моделей данные приводятся к единой шкале:
лог-доходности и волатильности годятся как есть, корреляции
преобразуются через преобразование Фишера $z = \operatorname{atanh}(\rho)$
для расширения диапазона входов. Для борьбы с тяжёлыми хвостами
лог-доходностей используется winsorization на уровне квантилей $0{.}5\%$
и $99{.}5\%$.

\todoblock{Если объём подраздела позволит, переместить эту главу
в приложение, оставив в основном тексте только сводку.}
